% !TEX root = ../Abschlussarbeit_Jan_Händl.tex
\appendix
\section{Anhang}

\subsection{Hilfsmittel}
\label{subsec:hilfsmittel}
Folgend werden die genutzten Hilfsmittel aufgelistet.
\begin{table}[H]
  \centering
  \renewcommand{\arraystretch}{1.4} % Mehr Zeilenabstand
  \setlength{\tabcolsep}{12pt}      % Mehr Spaltenabstand
  \caption{Hilfsmittel}
  \label{tab:hilfsmittel}
  \begin{tabularx}{\textwidth}{LLL}
    \toprule
    \textbf{Hilfsmittel}                                        & \textbf{Werkzeuge}                                                   & \textbf{Verwendung}                                      \\
    \midrule
    Zeichentool für Skizzen und Diagramme                       & Mermaid, Excalidraw, Screenshots                                     & Abb. 2.1, 2.2, 2.3, 2.4, 3.4, 4.1                        \\ \hline
    Rechtschreib- und Grammatikprüfung                          & Gemini 3.5, LTeX+ grammar/spell checking using LanguageTool  v15.7.1 & Komplette Arbeit                                         \\ \hline
    Schreibprogramm                                             & Visual Studio Code, LaTeX Workshop                                   & Komplette Arbeit                                         \\ \hline
    Version Control System                                      & Git                                                                  & Komplette Arbeit                                         \\ \hline
    Literaturrecherche                                          & Google Scholar, Researchgate, NotebookLM                             & Komplette Arbeit                                         \\ \hline
    Integrated Development Environment                          & Visual Studio Code                                                   & Programmierteil (siehe USB-Stick)                        \\ \hline
    Deployment Neo4j Datenbank                                  & Docker Compose                                                       & Programmierteil (siehe USB-Stick)                        \\ \hline
    Visualisierung Neo4j Graph                                  & Neo4j Browser                                                        & Abb. 3.1, 3.2, 3.4                                       \\ \hline
    Generierung Code und GraphQL Schema, Resolver und Mutations & Gemini 3.1, Gemini 3.5                                               & Programmierteil (siehe USB-Stick), Listing. 6, 8, 10, 11 \\ \hline
    Konfiguration Latex Projekt                                 & Gemini 3.1, Gemini 3.5                                               & Komplette Arbeit                                         \\

    \bottomrule
  \end{tabularx}
\end{table}

\subsection{Verwendete KI-Prompts}
\label{subsec:ki_prompts}
Folgend sind die maßgeblichen KI-Prompts für die Arbeit aufgelistet.
\begin{table}[H]
  \centering
  \renewcommand{\arraystretch}{1.4} % Mehr Zeilenabstand
  \setlength{\tabcolsep}{12pt}      % Mehr Spaltenabstand
  \caption{Verwendete KI-Prompts}
  \label{tab:ki_prompts}
  \begin{tabular}{lll}
    \toprule
    \textbf{Prompt}        & \textbf{Ergebnis}             & \textbf{KI Modell}     \\
    \midrule
    Formuliere den Text um & Vorschläge für Formulierungen & Gemini 3.5, Gemini 3.1 \\
    \bottomrule
  \end{tabular}
\end{table}

\subsection{Routing in Microfrontends}
\label{subsec:microfrontends_routing}
In diesem Abschnitt wird der exemplarische Vergleich zwischen dem URL-basierten Routing und einem zustandsbasierten Routing (SPA) dargestellt.
\begin{lstlisting}[language=Javascript, caption={Exemplarischer Vergleich zwischen Url Routing und einer SPA}, label={lst:microfrontends_routing}]
import React, { useState } from 'react';
import { BrowserRouter, Routes, Route, Link } from 'react-router-dom';
// 1. Komponenten für das State-basierte Routing (SPA)
const RoombookApp = React.lazy(() =>
  import("roombook/RoombookApp").catch(() => {
    return {
      default: () => (
        <div className="p-8 text-center text-destructive bg-destructive/10 rounded-xl m-8">
          Fehler beim Laden von Roombook! Ist die Remote App (Port 5174)
          gestartet?
        </div>
      ),
    };
  }),
);
const InsightsApp = React.lazy(() =>
  import("insights/InsightsApp").catch(() => {
    return {
      default: () => (
        <div className="p-8 text-center text-destructive bg-destructive/10 rounded-xl m-8">
          Fehler beim Laden von Insights! Ist die Remote App (Port 5175)
          gestartet?
        </div>
      ),
    };
  }),
);
// 2. Komponenten für das URL-basierte Routing
const RoombookApp = () => <p>Komponente: RoombookApp (URL)</p>;
const InsightsApp = () => <p>Komponente: InsightsApp (URL)</p>;
export default function RoutingVergleichApp() {
  // State für das zustandsbasierte Umschalten
  const [activeTab, setActiveTab] = useState('A');
  return (
    <BrowserRouter>
      {/* MECHANISMUS 1: URL-BASIERTES ROUTING (Ändert die Browser-URL)*/}
      <div>
        <h3>1. URL-basiertes Routing</h3>
        <nav>
          <Link to="/roombookapp"/>
          <Link to="/insightsapp"/>
        </nav>

        <Routes>
          <Route path="/roombookapp" element={<RoombookApp />} />
          <Route path="/insightsapp" element={<InsightsApp />} />
        </Routes>
      </div>
      <hr />
      {/*MECHANISMUS 2: STATE-BASIERTES ROUTING (SPA)*/}
      <div>
        <h3>2. State-basiertes Routing</h3>
        <nav>
          <button onClick={() => setActiveTab('A')}>Zeige RoombookApp</button>
          <button onClick={() => setActiveTab('B')}>Zeige InsightsApp</button>
        </nav>
        <div>
          {activeTab === 'A' && <RoombookApp />}
          {activeTab === 'B' && <InsightsApp />}
        </div>
      </div>
    </BrowserRouter>
  );
}
\end{lstlisting}

\subsection{Cypher-Datenbankabfragen}
\label{subsec:cypher_abfragen}
Folgend wird ein exemplarisches Cypher-Statement aufgeführt, welches zur Abfrage von Sensorbeziehungen in der Neo4j-Graphdatenbank verwendet wird.
\begin{lstlisting}[language=cypher, caption={Exemplarische Cypher-Abfrage für aktive Sensoren in einem Raum}, label={lst:cypher_active_sensors}]
MATCH (r:Room {number: "1.04"})<-[:LOCATED_IN]-(s:Sensor)
WHERE s.status = "ACTIVE"
RETURN r, s;
\end{lstlisting}

\subsection{Ergänzende Diagramme und Daten}
\label{subsec:ergaenzende_diagramme}
Hier können weitere unterstützende Diagramme wie UML-Klassendiagramme, Flussdiagramme oder zusätzliche Messdaten eingefügt werden, die den Rahmen des Hauptteils sprengen würden.
Beispielsweise zeigt ein Systemarchitektur-Diagramm die Interaktion zwischen den Microfrontends und dem API-Gateway über GraphQL.
\begin{center}
  \textit{[Hier kann ein zusätzliches Architekturdiagramm oder Flussdiagramm eingefügt werden]}
\end{center}